\begin{center}
  {\Huge \color{mainblue}\textbf{Esercitazione: volume e capacità}}\\[0.2cm]
  \rule{0.82\linewidth}{0.5pt}
\end{center}


\section*{Richiamo iniziale}

\begin{esercizio}{La merenda della classe}
Per una merenda in classe, l'insegnante porta alcune bevande:

\begin{itemize}
  \item una bottiglia grande di succo da $1{,}5\,\mathrm{l}$;
  \item $4$ bottigliette d'acqua da $500\,\mathrm{ml}$ ciascuna;
  \item $6$ lattine da $20\,\mathrm{cl}$ ciascuna.
\end{itemize}

\domanda{Quanti litri di bevande ci sono in tutto?}

\vspace{0.5cm}

\noindent
Scala delle misure di capacità:

\vspace{0.2cm}

\immagine[5cm]{immagini/scala_capacita}

\responsespace{5.8cm}

\begin{soluzione}
  \textbf{Passaggi:}
  \begin{itemize}
    \item \textbf{Succo:} $1{,}5\unita{l}$
    \item \textbf{Acqua:} $4 \times 500\unita{ml} = 2000\unita{ml} = 2\unita{l}$
    \item \textbf{Lattine:} $6 \times 20\unita{cl} = 120\unita{cl} = 1{,}2\unita{l}$
    \item \textbf{Totale:} $1{,}5\unita{l} + 2\unita{l} + 1{,}2\unita{l} = 4{,}7\unita{l}$
  \end{itemize}
\end{soluzione}
\end{esercizio}

\condnewpage

\section*{1. Problemi sulle misure di capacità}

\begin{esercizio}{Acqua per il campeggio}
Per un campeggio, una classe porta una cisterna da $1\,\mathrm{hl}$ di acqua potabile. Ogni studente riceve $2\,\mathrm{l}$ di acqua al giorno e il gruppo è composto da $10$ ragazzi.

\domanda{Per quanti giorni durerà la scorta?}

\responsespace{7cm}

\begin{soluzione}
  \textbf{Passaggi:}
  \begin{enumerate}
    \item Convertiamo la capacità della cisterna in litri: $1\unita{hl} = 100\unita{l}$.
    \item Calcoliamo il consumo totale giornaliero del gruppo: $10\text{ ragazzi} \times 2\unita{l/giorno} = 20\unita{l/giorno}$.
    \item Dividiamo la scorta totale per il consumo giornaliero: $\frac{100\unita{l}}{20\unita{l/giorno}} = 5\text{ giorni}$.
  \end{enumerate}
\end{soluzione}
\end{esercizio}

\begin{esercizio}{Il serbatoio della macchina del caffè}
Nel bar della scuola, una macchina del caffè ha un serbatoio che può contenere $6\,\mathrm{l}$ di acqua. Ogni caffè consuma $15\,\mathrm{cl}$ d'acqua.

\domanda{Quanti caffè si possono preparare prima di dover riempire di nuovo il serbatoio?}

\responsespace{5.8cm}

\begin{soluzione}
  \textbf{Passaggi:}
  \begin{enumerate}
    \item Convertiamo la capacità del serbatoio in centilitri per facilitare la divisione: $6\unita{l} = 600\unita{cl}$.
    \item Dividiamo il volume del serbatoio per il consumo di un singolo caffè: $\frac{600\unita{cl}}{15\unita{cl/caff\grave{e}}} = 40\text{ caffè}$.
  \end{enumerate}
\end{soluzione}
\end{esercizio}

\condnewpage

\section*{2. Collegare capacità e volume}

\nota{Le misure di capacità e le misure di volume sono collegate. La relazione più importante è
\[
1\,\mathrm{dm^3}=1\,\mathrm{l}.
\]
Questo significa che un cubo di lato $1\,\mathrm{dm}$ può contenere esattamente $1\,\mathrm{l}$ di liquido.}

\immagine[4.5cm]{immagini/relazione_litro_dm3}

\begin{esercizio}{Acqua nel contenitore graduato}
In laboratorio, una studentessa riempie un contenitore graduato con $12\,\mathrm{dl}$ d'acqua.

\begin{enumerate}[label=\alph*)]
  \item Qual è la capacità dell'acqua contenuta nel recipiente, espressa in litri?
  \item Qual è il volume dell'acqua contenuta nel recipiente, espresso in $\mathrm{dm^3}$?
  \item Qual esprimi il volume dell'acqua contenuta nel recipiente, espresso in $\mathrm{cm^3}$?
\end{enumerate}

Scala di misura del volume: 
\immagine[4.5cm]{immagini/add}

\responsespace{6.4cm}

\begin{soluzione}
  \textbf{Passaggi:}
  \begin{enumerate}[label=\alph*)]
    \item Convertiamo i decilitri in litri: $12\unita{dl} = 1{,}2\unita{l}$.
    \item Poiché $1\unita{l} = 1\unita{dm^3}$, allora $1{,}2\unita{l} = 1{,}2\unita{dm^3}$.
    \item Convertiamo i decimetri cubi in centimetri cubi. Sulla scala del volume facciamo un passo verso destra (moltiplicando per $1000$): $1{,}2\unita{dm^3} \times 1000 = 1200\unita{cm^3}$.
  \end{enumerate}
\end{soluzione}
\end{esercizio}

\condnewpage

\begin{esercizio}{Il latte per il dolce}
In cucina, un cuoco usa un misurino per versare $4{,}2\,\mathrm{dl}$ di latte in una ciotola per preparare un dolce.

\begin{enumerate}[label=\alph*)]
  \item Esprimi la quantità di latte in litri.
  \item Esprimi il volume del latte in $\mathrm{mm^3}$.
\end{enumerate}

\responsespace{10cm}

\begin{soluzione}
  \textbf{Passaggi:}
  \begin{enumerate}[label=\alph*)]
    \item Convertiamo decilitri in litri: $4{,}2\unita{dl} = 0{,}42\unita{l}$.
    \item Sappiamo che $0{,}42\unita{l} = 0{,}42\unita{dm^3}$. Per passare da decimetri cubi a millimetri cubi facciamo due salti verso destra sulla scala dei volumi ($\unita{dm^3} \to \unita{cm^3} \to \unita{mm^3}$), ovvero moltiplichiamo per $1\,000\,000$ ($1000 \times 1000$):
    \[
    0{,}42\unita{dm^3} \times 1\,000\,000 = 420\,000\unita{mm^3}
    \]
  \end{enumerate}
\end{soluzione}
\end{esercizio}

\begin{esercizio}{Uguaglianze da completare}
Completa le uguaglianze. Puoi usare la scala delle capacità e la relazione tra litro e decimetro cubo.

\vspace{0.2cm}
\renewcommand{\arraystretch}{2.25}
\begin{tabularx}{\linewidth}{|X|X|}
\hline
$1\,\mathrm{l}=\fillbox{3.1cm}\,\mathrm{dl}$ & $1\,\mathrm{l}=\fillbox{3.1cm}\,\mathrm{cl}$ \\
\hline
$1\,\mathrm{l}=\fillbox{3.1cm}\,\mathrm{ml}$ & $1\,\mathrm{dm^3}=\fillbox{3.1cm}\,\mathrm{l}$ \\
\hline
$1\,\mathrm{cm^3}=\fillbox{3.1cm}\,\mathrm{ml}$ & $1\,\mathrm{m^3}=\fillbox{3.1cm}\,\mathrm{l}$ \\
\hline
$3\,\mathrm{l}=\fillbox{3.1cm}\,\mathrm{ml}$ & $250\,\mathrm{cl}=\fillbox{3.1cm}\,\mathrm{l}$ \\
\hline
\end{tabularx}

\begin{soluzione}
  \textbf{Soluzione delle uguaglianze:}
  \begin{itemize}
    \item $1\unita{l} = \mathbf{10}\unita{dl}$ \hfill $1\unita{l} = \mathbf{100}\unita{cl}$
    \item $1\unita{l} = \mathbf{1000}\unita{ml}$ \hfill $1\unita{dm^3} = \mathbf{1}\unita{l}$
    \item $1\unita{cm^3} = \mathbf{1}\unita{ml}$ \hfill $1\unita{m^3} = \mathbf{1000}\unita{l}$
    \item $3\unita{l} = \mathbf{3000}\unita{ml}$ \hfill $250\unita{cl} = \mathbf{2{,}5}\unita{l}$
  \end{itemize}
\end{soluzione}
\end{esercizio}

\condnewpage

\section*{3. Recipienti a forma di parallelepipedo}

\begin{esercizio}{Un secchio a base quadrata}
Il seguente secchio ha per base un quadrato di lato $10\,\mathrm{cm}$ ed è alto $30\,\mathrm{cm}$.

\immagine[3.1cm]{immagini/\theesercizio}

\begin{enumerate}[label=\alph*)]
  \item Calcola il suo volume.
  \item Quanti litri d'acqua potrebbe contenere se venisse riempito fino all'orlo?
\end{enumerate}

\responsespace{4.2cm}

\begin{soluzione}
  \textbf{Passaggi:}
  \begin{enumerate}[label=\alph*)]
    \item Calcoliamo l'area della base quadrata: $A_b = 10\unita{cm} \times 10\unita{cm} = 100\unita{cm^2}$.\\
    Calcoliamo il volume: $V = A_b \times h = 100\unita{cm^2} \times 30\unita{cm} = 3000\unita{cm^3}$.
    \item Convertiamo i centimetri cubi in decimetri cubi dividendo per $1000$: $3000\unita{cm^3} = 3\unita{dm^3}$.\\
    Poiché $1\unita{dm^3} = 1\unita{l}$, il secchio contiene $3\unita{l}$.
  \end{enumerate}
\end{soluzione}
\end{esercizio}
\condnewpage
\begin{esercizio}{La piscina}
La piscina raffigurata è lunga $14\,\mathrm{m}$, larga $6\,\mathrm{m}$ e profonda $1{,}8\,\mathrm{m}$.

\domanda{Quanti litri d'acqua sono necessari per riempirla?}

\immagine[5cm]{immagini/\theesercizio}

\responsespace{4.5cm}

\begin{soluzione}
  \textbf{Passaggi:}
  \begin{enumerate}
    \item Calcoliamo il volume della piscina in metri cubi:
    \[
    V = 14\unita{m} \times 6\unita{m} \times 1{,}8\unita{m} = 151{,}2\unita{m^3}
    \]
    \item Convertiamo i metri cubi in decimetri cubi ($\unita{dm^3}$ ovvero litri) moltiplicando per $1000$:
    \[
    151{,}2\unita{m^3} \times 1000 = 151\,200\unita{dm^3} = 151\,200\unita{l}
    \]
  \end{enumerate}
\end{soluzione}
\end{esercizio}

\condnewpage

\begin{esercizio}{La vaschetta di vetro}
La seguente vaschetta di vetro, non piena, contiene $30\,\mathrm{l}$ di acqua. La base misura $4\,\mathrm{dm}$ per $3\,\mathrm{dm}$.

\domanda{Quale altezza raggiunge l'acqua al suo interno?}

\immagine[5cm]{immagini/\theesercizio}

\responsespace{5.8cm}
\condnewpage

\begin{soluzione}
  \textbf{Passaggi:}
  \begin{enumerate}
    \item Esprimiamo la quantità d'acqua come volume: $30\unita{l} = 30\unita{dm^3}$.
    \item Calcoliamo l'area della base della vaschetta: $A_b = 4\unita{dm} \times 3\unita{dm} = 12\unita{dm^2}$.
    \item Troviamo l'altezza dell'acqua in due modi:
    \begin{itemize}
      \item \textbf{1\textdegree\ modo (a tentativi):} L'area di base è $12\unita{dm^2}$. Cerchiamo l'altezza $h$ tale che:
      \[
      \text{Area base} \times \text{Altezza} = \text{Volume} \implies 12 \times h = 30
      \]
      Se proviamo $h = 2\unita{dm}$, il volume è $12 \times 2 = 24\unita{dm^3}$ (troppo piccolo).\\
      Se proviamo $h = 3\unita{dm}$, il volume è $12 \times 3 = 36\unita{dm^3}$ (troppo grande).\\
      Il valore corretto deve trovarsi esattamente a metà strada tra $2\unita{dm}$ e $3\unita{dm}$, cioè $h = 2{,}5\unita{dm}$ (infatti $12 \times 2{,}5 = 30\unita{dm^3}$).
      \item \textbf{2\textdegree\ modo (formula inversa):} Dividiamo il volume per l'area di base:
      \[
      h = \frac{V}{A_b} = \frac{30\unita{dm^3}}{12\unita{dm^2}} = 2{,}5\unita{dm} = 25\unita{cm}
      \]
    \end{itemize}
  \end{enumerate}
\end{soluzione}
\end{esercizio}

\begin{esercizio}{Una bacinella senza coperchio}
La figura seguente rappresenta una bacinella di latta priva di coperchio. Le sue dimensioni sono: lunghezza $50\,\mathrm{cm}$, larghezza $40\,\mathrm{cm}$, altezza $20\,\mathrm{cm}$.

\immagine[5cm]{immagini/\theesercizio}

\begin{enumerate}[label=\alph*)]
  \item Quanti $\mathrm{dm^2}$ di latta sono necessari per costruirla?
  \item Quanto misura il suo volume?
  \item Se venisse riempita completamente, quanti litri d'acqua potrebbe contenere?
\end{enumerate}

\responsespace{4.2cm}

\begin{soluzione}
  \textbf{Passaggi:}
  \begin{enumerate}[label=\alph*)]
    \item Convertiamo le dimensioni in decimetri: lunghezza $5\unita{dm}$, larghezza $4\unita{dm}$, altezza $2\unita{dm}$.\\
    Calcoliamo l'area delle pareti esterne (senza coperchio):
    \begin{itemize}
      \item \textbf{Base (fondo):} $5\unita{dm} \times 4\unita{dm} = 20\unita{dm^2}$
      \item \textbf{2 pareti lunghe:} $2 \times (5\unita{dm} \times 2\unita{dm}) = 20\unita{dm^2}$
      \item \textbf{2 pareti corte:} $2 \times (4\unita{dm} \times 2\unita{dm}) = 16\unita{dm^2}$
      \item \textbf{Totale latta:} $20 + 20 + 16 = 56\unita{dm^2}$.
    \end{itemize}
    \item Calcoliamo il volume: $V = 5\unita{dm} \times 4\unita{dm} \times 2\unita{dm} = 40\unita{dm^3}$ (pari a $40\,000\unita{cm^3}$).
    \item Poiché $1\unita{dm^3} = 1\unita{l}$, può contenere esattamente $40\unita{l}$.
  \end{enumerate}
\end{soluzione}
\end{esercizio}

\condnewpage

\section*{4. Problemi di applicazione}

\begin{esercizio}{La vasca per pesci}
Per riempire una vasca per pesci fino a $3\,\mathrm{cm}$ dal bordo superiore bisogna usare $6$ volte il misurino da $2\,\mathrm{l}$. La base della vasca misura $30\,\mathrm{cm}$ per $16\,\mathrm{cm}$.

\domanda{Calcola l'altezza della vasca.}

\immagine[6cm]{immagini/\theesercizio}


\responsespace{6.2cm}
\condnewpage

\begin{soluzione}
  \textbf{Passaggi:}
  \begin{enumerate}
    \item Calcoliamo il volume di acqua inserito: $6 \times 2\unita{l} = 12\unita{l} = 12\unita{dm^3} = 12\,000\unita{cm^3}$.
    \item Calcoliamo l'area della base in $\unita{cm^2}$: $A_b = 30\unita{cm} \times 16\unita{cm} = 480\unita{cm^2}$.
    \item Troviamo l'altezza raggiunta dall'acqua ($h_{\text{acqua}}$) in due modi:
    \begin{itemize}
      \item \textbf{1\textdegree\ modo (a tentativi):} L'area di base è $480\unita{cm^2}$. Cerchiamo l'altezza $h$ tale che:
      \[
      \text{Area base} \times \text{Altezza} = \text{Volume} \implies 480 \times h = 12\,000
      \]
      Se proviamo $h = 10\unita{cm}$, il volume è $480 \times 10 = 4800\unita{cm^3}$ (troppo piccolo).\\
      Se proviamo $h = 20\unita{cm}$, il volume è $480 \times 20 = 9600\unita{cm^3}$ (troppo piccolo).\\
      Se proviamo $h = 30\unita{cm}$, il volume è $480 \times 30 = 14\,400\unita{cm^3}$ (troppo grande).\\
      Proviamo a metà strada tra $20\unita{cm}$ e $30\unita{cm}$, cioè $h = 25\unita{cm}$: infatti $480 \times 25 = 12\,000\unita{cm^3}$ (perfetto!).
      \item \textbf{2\textdegree\ modo (formula inversa):} Dividiamo il volume dell'acqua per l'area di base:
      \[
      h_{\text{acqua}} = \frac{V}{A_b} = \frac{12\,000\unita{cm^3}}{480\unita{cm^2}} = 25\unita{cm}
      \]
    \end{itemize}
    \item Aggiungiamo lo spazio vuoto di $3\unita{cm}$ per trovare l'altezza totale: $h_{\text{vasca}} = 25\unita{cm} + 3\unita{cm} = 28\unita{cm}$.
  \end{enumerate}
\end{soluzione}
\end{esercizio}

\begin{esercizio}{Lunghezza della piscina}
Questa piscina è profonda $2\,\mathrm{m}$ ed è larga $8\,\mathrm{m}$. Per riempirla occorrono $240\,000\,\mathrm{l}$.

\domanda{Calcola la lunghezza della piscina.}

\immagine[5cm]{immagini/\theesercizio}

\responsespace{4.8cm}

\begin{soluzione}
  \textbf{Passaggi:}
  \begin{enumerate}
    \item Convertiamo la capacità dell'acqua in volume ($\unita{m^3}$): $240\,000\unita{l} = 240\,000\unita{dm^3} = 240\unita{m^3}$.
    \item Troviamo l'area della sezione verticale (larghezza $\times$ profondità): $A = 8\unita{m} \times 2\unita{m} = 16\unita{m^2}$.
    \item Troviamo la lunghezza della piscina ($l$) in due modi:
    \begin{itemize}
      \item \textbf{1\textdegree\ modo (a tentativi):} Sappiamo che la sezione è $16\unita{m^2}$. Cerchiamo la lunghezza $l$ tale che:
      \[
      \text{Sezione} \times \text{Lunghezza} = \text{Volume} \implies 16 \times l = 240
      \]
      Se proviamo $l = 10\unita{m}$, il volume è $16 \times 10 = 160\unita{m^3}$ (troppo piccolo).\\
      Se proviamo $l = 20\unita{m}$, il volume è $16 \times 20 = 320\unita{m^3}$ (troppo grande).\\
      Proviamo a metà strada tra $10\unita{m}$ e $20\unita{m}$, cioè $l = 15\unita{m}$: infatti $16 \times 15 = 240\unita{m^3}$ (perfetto!).
      \item \textbf{2\textdegree\ modo (formula inversa):} Dividiamo il volume per l'area della sezione:
      \[
      l = \frac{V}{\text{larghezza} \times \text{profondità}} = \frac{240\unita{m^3}}{16\unita{m^2}} = 15\unita{m}
      \]
    \end{itemize}
  \end{enumerate}
\end{soluzione}
\end{esercizio}

\condnewpage

\begin{esercizio}{Tempo di riempimento}
Una vasca per pesci ha la forma di un parallelepipedo rettangolo. Le dimensioni di base sono $5\,\mathrm{m}$ e $3\,\mathrm{m}$, con un'altezza di $1{,}2\,\mathrm{m}$.

Per riempirla vengono utilizzati due rubinetti, da cui fuoriescono complessivamente $9\,\mathrm{l}$ di acqua al minuto.

\domanda{Calcola quanto tempo sarà necessario per riempire la vasca.}

\immagine[5cm]{immagini/\theesercizio}

\responsespace{7cm}

\begin{soluzione}
  \textbf{Passaggi:}
  \begin{enumerate}
    \item Calcoliamo il volume della vasca in $\unita{m^3}$: $V = 5\unita{m} \times 3\unita{m} \times 1{,}2\unita{m} = 18\unita{m^3}$.
    \item Convertiamo il volume in litri: $18\unita{m^3} = 18\,000\unita{l}$.
    \item Dividiamo il volume totale per l'acqua che esce complessivamente al minuto dai rubinetti per ottenere il tempo in minuti:
    \[
    t = \frac{18\,000\unita{l}}{9\unita{l/min}} = 2000\text{ minuti}
    \]
    \item (Opzionale) Esprimiamo il tempo in ore dividendo per $60$:
    \[
    2000\text{ min} = 33\text{ ore e } 20\text{ minuti} \quad (\text{poiché } 33 \times 60 = 1980\text{ min})
    \]
  \end{enumerate}
\end{soluzione}
\end{esercizio}

\condnewpage

\begin{esercizio}{Due cartoni di tè freddo}
Antonio e Bea sono assetati. Entrambi hanno con sé un cartone di tè freddo. Supponiamo, per semplificare, che nei due contenitori non ci sia aria, ma solo tè.

Bea ha più tè di Antonio e ne offre un po' ad Antonio. Alla fine, quando i due contenitori sono vuoti, entrambi hanno bevuto la stessa quantità di liquido.

\immagine[5cm]{immagini/\theesercizio}

\begin{enumerate}[label=\roman*)]
  \item Quanto tè ha offerto Bea ad Antonio?
  \item Quale dei due contenitori richiede più cartone per la fabbricazione? Per semplicità non si tiene conto delle linguette e delle pieghe.
\end{enumerate}

\responsespace{2.8cm}

\begin{soluzione}
  \textbf{Passaggi:}
  \begin{enumerate}[label=\roman*)]
    \item \textbf{Volume contenitore Antonio ($9{,}6\unita{cm} \times 8\unita{cm} \times 15\unita{cm}$):}\\
    $V_A = 9{,}6 \times 8 \times 15 = 1152\unita{cm^3} = 1{,}152\unita{l}$.\\
    \textbf{Volume contenitore Bea ($12{,}5\unita{cm} \times 6{,}4\unita{cm} \times 8\unita{cm}$):}\\
    $V_B = 12{,}5 \times 6{,}4 \times 8 = 640\unita{cm^3} = 0{,}640\unita{l}$.\\
    \textbf{Volume totale di tè:} $1152 + 640 = 1792\unita{cm^3}$.\\
    Ognuno deve bere la metà del totale: $\frac{1792}{2} = 896\unita{cm^3} = 896\unita{ml}$.\\
    Bea offre ad Antonio la quantità necessaria a portarlo a $896\unita{ml}$: $896 - 640 = 256\unita{ml}$ (o $0{,}256\unita{l}$).
    \item Calcoliamo la superficie totale dei due cartoni:
    \begin{itemize}
      \item \textbf{Superficie Antonio:} $S_A = 2 \times (9{,}6 \times 8 + 9{,}6 \times 15 + 8 \times 15) = 2 \times (76{,}8 + 144 + 120) = 681{,}6\unita{cm^2}$.
      \item \textbf{Superficie Bea:} $S_B = 2 \times (12{,}5 \times 6{,}4 + 12{,}5 \times 8 + 6{,}4 \times 8) = 2 \times (80 + 100 + 51{,}2) = 462{,}4\unita{cm^2}$.
    \end{itemize}
    Il contenitore di Antonio richiede più cartone.
  \end{enumerate}
\end{soluzione}
\end{esercizio}

\condnewpage

\section*{5. Sfide finali}

\begin{esercizio}{Profondità di una piscina}
Una piscina a forma di parallelepipedo rettangolo viene riempita con sei rubinetti, da cui fuoriescono complessivamente $12\,\mathrm{l}$ di acqua al minuto. Le dimensioni della base della piscina sono $25\,\mathrm{m}$ e $15\,\mathrm{m}$.

Per riempirla completamente è stato necessario un tempo di $6$ giorni, $12$ ore e $15$ minuti.

\domanda{Calcola la profondità della piscina.}

\responsespace{8cm}

\begin{soluzione}
  \textbf{Passaggi:}
  \begin{enumerate}
    \item Calcoliamo il tempo totale di riempimento in minuti:
    \begin{itemize}
      \item $6\text{ giorni} = 6 \times 24 \times 60 = 8640\text{ min}$
      \item $12\text{ ore} = 12 \times 60 = 720\text{ min}$
      \item $15\text{ minuti} = 15\text{ min}$
      \item \textbf{Tempo totale:} $8640 + 720 + 15 = 9375\text{ minuti}$.
    \end{itemize}
    \item Calcoliamo il volume di acqua immesso totale:
    \[
    V = 9375\unita{min} \times 12\unita{l/min} = 112\,500\unita{l} = 112\,500\unita{dm^3} = 112{,}5\unita{m^3}
    \]
    \item Calcoliamo l'area della base della piscina: $A_b = 25\unita{m} \times 15\unita{m} = 375\unita{m^2}$.
    \item Troviamo la profondità ($h$) in due modi:
    \begin{itemize}
      \item \textbf{1\textdegree\ modo (a tentativi):} L'area di base è $375\unita{m^2}$. Cerchiamo la profondità $h$ tale che:
      \[
      \text{Area base} \times \text{Profondità} = \text{Volume} \implies 375 \times h = 112{,}5
      \]
      Se proviamo $h = 1\unita{m}$, il volume è $375 \times 1 = 375\unita{m^3}$ (troppo grande).\\
      Se proviamo $h = 0{,}1\unita{m}$ ($10\unita{cm}$), il volume è $375 \times 0{,}1 = 37{,}5\unita{m^3}$ (troppo piccolo).\\
      Se proviamo $h = 0{,}5\unita{m}$ ($50\unita{cm}$), il volume è $375 \times 0{,}5 = 187{,}5\unita{m^3}$ (troppo grande).\\
      Se proviamo $h = 0{,}3\unita{m}$ ($30\unita{cm}$), il volume è $375 \times 0{,}3 = 112{,}5\unita{m^3}$ (perfetto!).
      \item \textbf{2\textdegree\ modo (formula inversa):} Dividiamo il volume totale per l'area di base:
      \[
      h = \frac{V}{A_b} = \frac{112{,}5\unita{m^3}}{375\unita{m^2}} = 0{,}3\unita{m} = 30\unita{cm}
      \]
    \end{itemize}
  \end{enumerate}
\end{soluzione}
\end{esercizio}
\condnewpage
\begin{esercizio}{Oggetto immerso}
I due recipienti raffigurati contengono la stessa quantità di acqua. Nel secondo recipiente è stato immerso un oggetto.

\immagine[6cm]{immagini/\theesercizio}

\domanda{Calcola il volume dell'oggetto immerso nel secondo recipiente.}

\responsespace{8cm}

\begin{soluzione}
  \textbf{Passaggi:}
  \begin{enumerate}
    \item Calcoliamo l'area di base del recipiente: $A_b = 20\unita{cm} \times 10\unita{cm} = 200\unita{cm^2}$.
    \item Calcoliamo l'innalzamento del livello dell'acqua dovuto all'immersione dell'oggetto:
    \[
    \Delta h = 13\unita{cm} - 10\unita{cm} = 3\unita{cm}
    \]
    \item Il volume dell'oggetto immerso corrisponde al volume dell'acqua spostata:
    \[
    V_{\text{oggetto}} = A_b \times \Delta h = 200\unita{cm^2} \times 3\unita{cm} = 600\unita{cm^3}
    \]
    \item (Opzionale) Esprimiamo il valore anche in decimetri cubi e litri per completezza: $600\unita{cm^3} = 0{,}6\unita{dm^3} = 0{,}6\unita{l}$.
  \end{enumerate}
\end{soluzione}
\end{esercizio}

\end{document}
